This paper addressed blink detection in the epoch-structured setting, a practically important but largely overlooked regime in which a continuous recording is divided into non-overlapping analysis windows before automated processing begins. We argued that the standard practice of concatenating available epochs and deriving a single amplitude threshold from the pooled signal can be problematic, because blink-free epochs may dilute the amplitude statistics that govern threshold placement. To address this, we proposed a three-stage epoch-aware pipeline: Stage~A screens epochs by their peak-to-peak amplitude using an autoreject-inspired criterion, Stage~B estimates a sample-level threshold from the resulting pool of suspicious epochs using the median and median absolute deviation, and Stage~C extracts blink regions and maps them back to epoch-local coordinates, with a fallback that stabilises estimation when too few suspicious epochs are available. We evaluated four conditions across two naturalistic driving corpora under an event-level overlap criterion with macro-averaged $F_1$ and Wilcoxon signed-rank tests. The central findings are threefold. First, the proposed median configuration improved over the BLINKER-concat and MNE-annotation baselines, primarily by recovering precision without collapsing recall. Second, the robust MAD-based estimator (Proposed-Med) was preferable to its mean-based counterpart (Proposed-Mean), consistent with the expectation that the median and MAD reduce sensitivity to within-epoch peaks in the suspicious-epoch pool. Third, Proposed-Med was the best-performing non-baseline method and significantly outperformed BLINKER-concat, MNE-annot, and Proposed-Mean under Wilcoxon signed-rank tests with Bonferroni correction, while performance was relatively stable across the tested epoch durations and showed only a small Raja--Cao2018 macro-$F_1$ gap. We also report an honest limitation: under the selected 30-second epoch setting, the proposed median configuration did not fully resolve all threshold-placement challenges. Characterising precisely when Stage~A helps and when it is overly conservative is a clear direction for future work. Further avenues include extending the evaluation to synchronized EEG--video benchmarks with high-confidence blink-region labels, and studying threshold-family selection as a first-class problem under a shared protocol rather than as a choice absorbed into a larger pipeline.
